\section{Introduction}

Language models are increasingly deployed within \emph{harnesses}: the programs that specify their prompts, tools, control flow, context and memory, and the orchestration code invoking them
\citep{weng2023agent,wang2024survey,weng2026harness}. The same model can exhibit
substantially different capabilities under different harnesses
\citep{yao2026harnessbenchmeasuringharnesseffects, zhang2026bindingconstraint}, which makes \emph{harness
optimization} -- the iterative improvement of a harness on a
measured outcome under a fixed budget -- an increasingly important part of building capable AI systems. 

Recent work asks whether AI systems can automate this process, extending a
broader line on self-improving agents, agent meta-optimization, and automated
research
\citep{rank2026posttrainbenchllmagentsautomate,
meng2026rsibenchdatabenchmarkingdatacentricresearch,
wijk2025rebenchevaluatingfrontierai}. These approaches differ in the role they
assign the model. ShinkaEvolve
\citep{lange2025shinkaevolveopenendedsampleefficientprogram} and GEPA
\citep{agrawal2026gepa} use it as a mutation operator inside a larger
evolutionary procedure, whereas \vero\ \citep{ursekar2026vero} and MetaHarness
\citep{lee2026metaharnessendtoendoptimizationmodel} use coding agents as
end-to-end optimizers that edit the harness as a codebase.

% Declared one unit earlier than it is discussed, on purpose. A full-width float
% cannot be placed on the page it is declared on -- LaTeX's output routine only
% considers double floats at a page boundary -- so a figure* declared partway down
% page N is deferred to the top of page N+1. Moving the declaration back is the only
% way to pull the figure forward; stfloats/dblfloatfix add bottom placement on the
% later page and do not help with this.
\begin{figure*}[t]
\centering
\includegraphics[width=\linewidth]{figures/fig1_architecture.pdf}
\caption{\textbf{Trusted execution for held-out harness optimization.} The optimizer can write only the target agent's harness. It can read the evaluation results and the target task data but not modify them. The optimizer
receives per-case development traces and aggregate validation metrics, while test cases and scores remain inaccessible behind the trusted evaluation server. Each candidate is evaluated in isolated, unprivileged sandboxes, and every model call, both the optimizer's own and the evaluations', passes through a gateway that enforces model allow-lists and per-scope budgets. The test partition is evaluated only after the optimizer nominates a final candidate. Because held-out data, provider credentials, and budget enforcement are absent from the optimizer's sandbox, these restrictions are properties of the execution environment rather than instructions the optimizer is expected to follow.}
\label{fig:arch}
\end{figure*}

Which role a domain calls for depends on the cost of one reliable evaluation,
and that cost is what makes harness optimization a test of more than coding
ability. A test suite reports cheaply whether a code change is correct; the
effect of a harness change must be estimated by running a stochastic agent over
many cases, at substantial cost. An optimizer must therefore diagnose
failures from incomplete evidence, implement system-level changes, spend a
limited evaluation budget, separate real improvement from noise, and decide what
to deploy. Where evaluation is cheap, selection over many candidates substitutes
for reasoning about any one of them; where it is expensive and noisy, reasoning
from prior evidence becomes the capability itself. As the systems being
optimized grow more complex, evaluating them only gets more expensive, so the
second regime is the one that increasingly matters. Our
benchmark specifically targets tasks whose evaluation is itself costly and noisy (fixed-corpus
research, terminal use) rather than ones that are cheap to score.

Independent of its place in that self-improvement loop, harness optimization
is also long-horizon, plays out over a diverse and complex tool ecosystem
rather than a narrow action space, and requires reasoning-driven
interpretation of a stochastic system in pursuit of a measured metric rather
than production of a single correct step. Each property is already the
target of separate benchmarks; a task that combines all three is a demanding
test of frontier capability in its own right.

Measuring the capability requires controlling how apparent improvement can
arise. Methods are today evaluated with their own target agents, seeds, budgets,
disclosure policies, and scoring protocols, so their results conflate the
optimizer model, the coding harness it acts through, the target agent, and the
protocol. Separating them requires three conditions: the target model,
environment, and verifier held fixed; the final evaluation held out throughout
search, so improvement reflects generalization rather than fit to a visible
score; and a trusted execution boundary that enforces the budget, blocks access
to held-out state, and preserves every candidate for audit. Under these
conditions three questions become answerable: whether frontier models can be
distinguished at this task (\textbf{RQ1}), where they fall short
(\textbf{RQ2}), and how much the optimizer's own coding harness contributes
relative to the model (\textbf{RQ3}).

We make the following contributions:

\begin{itemize}[leftmargin=*,itemsep=2pt,topsep=3pt]

    \item \textbf{A controlled benchmark for harness optimization.}
    We introduce \heb, in which an optimizer---an LLM paired with a
    coding harness---receives a target agent's seed harness, graded
    development and validation feedback, and a fixed target-evaluation
    budget. It edits the harness and nominates a final candidate, which
    is scored by its normalized gain over the seed on a held-out test
    partition that remains inaccessible throughout search.

    \item \textbf{A trusted, reproducible evaluation protocol.}
    The suite comprises \NumSuiteTasks\ downstream tasks with pinned
    seeds, fixed and non-overlapping development, validation, and test
    splits, graded disclosure, and recorded baselines for both the seed
    and off-the-shelf harnesses. A trusted execution environment,
    building on \vero\ \citep{ursekar2026vero}, enforces access and
    target-evaluation budgets, isolates held-out state, meters resource
    use, and versions every candidate for audit
    (Figure~\ref{fig:arch}).

    \item \textbf{A controlled evaluation of optimizer models and coding
    harnesses.}
    We evaluate five frontier optimizer models under a shared coding
    harness and their respective native harnesses across
    \NumPaperTasks\ downstream tasks, yielding \NumScoredCells\ scored
    optimizer runs, and evaluate two additional coding harnesses on one
    task. Under this paired design, differences among optimizer models
    using the shared harness are larger on average than the corresponding
    differences between shared and native harnesses. Native harnesses
    provide no consistent advantage, and achievable gains vary
    substantially across tasks and seed regimes (\textbf{RQ1},
    \textbf{RQ3}).

    \item \textbf{Evidence of remaining capability gaps.}
    Release-level experiments show that \heb\ resolves variation among
    successive optimizer-model releases on OfficeQA
    (Section~\ref{sec:ladder}). Instrumented search trajectories show
    that broader intervention is associated with greater held-out gain,
    whereas detailed failure-trace inspection is rarely used and is not
    positively associated with gain
    (Section~\ref{sec:process}).

\end{itemize}

Although our experiments use coding agents as end-to-end optimizers,
\heb\ is agnostic to optimizer design: any system that operates within
the prescribed budget, edits the seed harness, and nominates a final
candidate can enter. In this way, \heb\ turns harness engineering from
an optimizer-specific demonstration into a reproducible evaluation
target for measuring and developing systems that improve AI agents.

% We introduce \heb, which realizes these conditions. An \emph{optimizer}---a
% model paired with a coding harness---receives a seed agent's code, development
% and validation evaluations, and a fixed evaluation budget; it edits the agent and
% nominates
% a final candidate, scored by its normalized gain over the seed on a test
% partition that is inaccessible during search. The suite comprises
% \NumSuiteTasks\ downstream tasks with pinned seeds, immutable $20/40/40$ splits,
% graded disclosure, and recorded baselines for both the seed and open-source
% harnesses (Table~\ref{tab:suite}), executing on a trusted harness that versions
% candidates and meters spend, built on \vero\ \citep{ursekar2026vero}.

% Using \heb, we evaluate five frontier models under a shared coding harness and
% under their native ones, over \NumScoredCells\ scored cells on \NumPaperTasks\
% tasks, plus two further coding harnesses on one task. The optimizer model
% matters roughly $\AxisRatio\times$ more than the harness it acts through, native
% harnesses are not consistently superior, and gain varies sharply by task and
% seed (RQ1, RQ3). Varying only the optimizer model across successive releases of
% two model families, one series tracks release order and the other does not, so the
% benchmark
% tracks model progress rather than saturating against it
% (Section~\ref{sec:ladder}). Instrumenting the optimizers' own trajectories and
% the harnesses they produce, breadth of intervention tracks gain across tasks
% whereas the volume of failure-trace inspection does not---and the deepest
% inspection the interface offers was invoked \NSpanCalls\ times across
% \NumScoredCells\ runs (RQ2,
% Section~\ref{sec:process}). \heb\ is agnostic to optimizer design: any program
% that works within the budget, edits the seed, and nominates a candidate can
% enter.

%%%%%%%%%%%%%%%%%%%%%%%%%%%%%%%%%%%%%%%%%%%
%%%%%%%%%%%%%%%%%%%%%%%%%%%%%%%%%%%%%%%%%%%
% Pre-revision version of the paragraphs above, superseded by the pass that
% reordered the introduction into premise / prior work / why the capability is
% hard / why measuring it needs controls / what we built. Kept for reference.
%
% [P1, objective clause only] ... We study \emph{harness optimization}, its iterative and
% evaluation-guided form: repeatedly modifying a harness to maximize one or more metrics under a fixed budget and operational constraints
% \citep{weng2026harness}.
%
% [P2] Recent work has begun to ask whether AI systems can automate this process as
% part of the broader goal of accelerating AI development
% \citep{rank2026posttrainbenchllmagentsautomate,
% meng2026rsibenchdatabenchmarkingdatacentricresearch,
% wijk2025rebenchevaluatingfrontierai}. Existing approaches optimize prompts,
% workflows, skills, context artifacts, or general harness code
% \citep{zhang2026dgm,zhang2026hyperagents,ursekar2026vero,
% lee2026metaharnessendtoendoptimizationmodel,zelikman2024stop,
% hu2025adas,robeyns2025sica,ye2026metacontextengineeringagentic}.
% Some use LLMs as mutation operators within a larger evolutionary procedure,
% as in ShinkaEvolve and GEPA
% \citep{lange2025shinkaevolveopenendedsampleefficientprogram,
% agrawal2026gepa}; others, including \vero\ and MetaHarness, use coding agents
% as end-to-end optimizers that edit a harness represented as
% a codebase
% \citep{ursekar2026vero,lee2026metaharnessendtoendoptimizationmodel}.
%
% [P3, the literature-gap paragraph and its five questions]
% This literature establishes harness optimization as a promising research
% direction, but it does not yet provide a common basis for measuring the
% underlying capability. Individual methods are generally evaluated using their
% own target agents, seed implementations, search budgets, information
% disclosure policies, and scoring protocols. Their reported results therefore
% conflate several distinct factors: the optimizer model, the coding harness
% through which it acts, the target agent being improved, and the protocol used
% to evaluate it. As a result, basic questions remain unanswered. Is harness
% optimization a capability that transfers across downstream tasks? How much of
% an optimizer's performance comes from the model rather than its coding
% harness? Do models perform best in their native scaffolds? How stable are
% results when an optimizer is run again? Most importantly, does a harness
% optimization benchmark improve as the optimizer model itself becomes more
% capable?
%
% [P4] Harness optimization is a particularly meaningful capability to measure
% because coding ability is necessary but insufficient. An effective optimizer
% must evaluate a stochastic system, diagnose failures from incomplete evidence,
% formulate and implement system-level changes, allocate a limited evaluation
% budget, distinguish real improvements from noise, and ultimately decide which
% candidate to deploy. Unlike conventional software-engineering tasks, it has no
% cheap and exact oracle available during development. A test suite can usually
% report immediately whether a code change is correct; the effect of a harness
% change must instead be estimated by repeatedly running a stochastic agent over
% multiple cases, often at substantial cost
% \citep{song2026paceproxyagenticcapability}. When such evaluations are cheap,
% selection over many candidates can compensate for limited reasoning about any
% single one. When they are expensive and noisy, each experiment becomes
% valuable, and the optimizer's ability to reason from prior evidence becomes
% part of the capability being measured.
%
% [P5] Turning this task into a valid benchmark requires controlling the channels
% through which apparent improvement can arise. First, the target model, target
% environment, and verifier must remain fixed, so that performance differences
% are attributable to changes in the harness. Second, the final evaluation must
% be held out throughout search, so that improvement reflects generalization
% rather than adaptation to a visible test score. Third, a trusted execution
% boundary must enforce search and evaluation budgets, prevent access to
% held-out state, constrain the target model and environment, and preserve every
% candidate for audit. Together, these conditions turn an open-ended engineering
% process into a controlled capability evaluation with a verifiable held-out
% measurement.
%
% [Contributions, pre-revision]
% \paragraph{Contributions.}
% We make four contributions:
%
% \begin{enumerate}[itemsep=1pt,topsep=2pt,leftmargin=*]
%   \item \textbf{A controlled benchmark suite.}
%   We introduce six downstream harness-optimization tasks spanning document QA,
%   fixed-corpus deep research, codebase QA, tool-agent dialogue, multi-step
%   reasoning, and terminal task solving. Each task contains a pinned seed agent,
%   an immutable $20/40/40$ development--validation--test split, and graded
%   information disclosure. Tasks execute within a trusted outer harness that
%   versions candidates, enforces per-scope budgets, and isolates held-out state
%   (Section~\ref{sec:benchmark}). The execution infrastructure builds on \vero
%   \citep{ursekar2026vero}.
%
%   \item \textbf{A factorial evaluation of optimizer models and coding
%   harnesses.}
%   We evaluate five optimizer models using both a shared, model-agnostic coding
%   harness and each model's native harness, yielding \NumScoredCells\ scored
%   cells across \NumPaperTasks\ tasks with two independent rounds per
%   configuration (Sections~\ref{sec:setup} and \ref{sec:results}). Changing the
%   optimizer model affects gain roughly $\AxisRatio\times$ more than changing
%   its coding harness; native harnesses are not consistently superior;
%   achievable improvement is strongly task-dependent; and rerun variability is
%   comparable to the gaps between adjacent systems. These findings support
%   reporting performance in tiers rather than as a fine-grained ranking.
%
%   \item \textbf{Evidence that the benchmark tracks model capability.}
%   Holding the task, target model, seed agent, coding harness, and budget fixed,
%   we vary only the optimizer model across \NumLadderRungs\ releases of one
%   family. Gain rises from \LadderLo\ to \LadderHi\
%   (Section~\ref{sec:ladder}), showing that the benchmark responds to advances
%   in the optimizer model rather than measuring only its scaffold or evaluation
%   noise.
%
%   \item \textbf{An analysis of optimizer behavior.}
%   We instrument the optimizers' trajectories to study which aspects of their
%   behavior predict held-out improvement (Section~\ref{sec:process}). The
%   breadth of harness components an optimizer modifies predicts gain across
%   tasks, whereas the volume of failure-trace inspection does not. The
%   interface's most detailed trace-inspection affordance is used only once in
%   the entire study, revealing a substantial gap between the optimization
%   process the benchmark enables and the one current systems actually perform.
% \end{enumerate}
%
% Although our experiments focus on coding agents as end-to-end optimizers,
% \heb\ is agnostic to optimizer design: any system that operates within the
% budget, edits the seed harness, and nominates a candidate can enter the
% benchmark. More broadly, \heb\ turns harness engineering from an
% optimizer-specific demonstration into a reproducible evaluation target,
% enabling the community to compare, diagnose, and develop systems that improve
% AI agents.
%%%%%%%%%%%%%%%%%%%%%%%%%%%%%%%%%%%%%%%%%%%
%%%%%%%%%%%%%%%%%%%%%%%%%%%%%%%%%%%%%%%%%%%

% Figure 1 moved to sections/task.tex: it is referenced only from Sections 3-4,
% and a full-width float on page 1 was the largest single consumer of
% introduction space.


